\section{Residual characters and profile-descent endpoints}
\label{app:li-comparison}

This appendix compares two predicates on balanced sextuples; it does not
assume an unrestricted classification from another preprint. Li's
September manuscript \cite[Theorem~1.1]{Li26a} already states the conclusion
of \Cref{thm:oddfour}. His November manuscript \cite[Theorem~A]{Li26b}
states the more specific positive classification described below.

\subsection{The predicates and their behaviour under inflation}
For an odd $m$, write $\mathrm{Dec}_m$ for the following sextuples:
those containing an inverse pair, and those admitting a partition into
triples $B,C$ with $\sum B=k\ne0$, $\sum C=-k$, such that
$B+\{-k\}$ and $\{k\}+C$ are balanced quadruples. This is the
surface-contraction criterion. Write $\mathrm{Std}_{5,m}$ for one
exact multiset $\sigma_{5,x}$ with $5x\ne0$, rather than an integral
combination of standard multisets. These are the predicates of
\cite[\S2.3]{Li26b}.

\begin{lemma}\label{lem:li-inflation}
Let $m=gd$ be odd and let $A$ be a sextuple of nonzero residues modulo
$d$ with sum zero. Multiplication by $g$ preserves and reflects
balance, membership in $\mathrm{Dec}$, and membership in
$\mathrm{Std}_5$. The primitive conductor of $gA$ is the primitive
conductor of $A$.
\end{lemma}
\begin{proof}
Every unit modulo $d$ lifts to a unit modulo $m$. Indeed, for each
prime dividing $m$ but not $d$, choose a nonzero residue in that prime
component and apply the Chinese remainder theorem, while preserving
the chosen unit modulo $d$. Also
$\overline{tga}_{m}=g\overline{ta}_{d}$. These statements prove the
balance assertion. Inverse pairs are preserved and reflected.
For a surface contraction of $gA$, the distinguished entry is the sum
of three entries of $gA$, hence is divisible by $g$; divide the two
balanced quadruples by $g$ and use the balance assertion. The converse
is obtained by multiplying a contraction by $g$.

If $gA$ is exactly $5$-standard, all five entries of its progression
are divisible by $g$. Their successive difference $m/5$ is therefore
divisible by $g$, so $5\mid d$. Its initial entry is $gx$ for a residue
$x$ modulo $d$; dividing the entire standard multiset by $g$ gives
$\sigma_{5,x}$ at level $d$. The condition $5gx\not\equiv0\pmod m$
is equivalent to $5x\not\equiv0\pmod d$. Inflation gives the converse.
Finally, the additive order of $ga$ modulo $gd$ is the order of $a$
modulo $d$. Taking the least common multiple proves the last assertion.
\end{proof}

\begin{figure}[t]
\centering
\begin{tikzpicture}[>=Stealth,font=\small,
 b/.style={draw,rounded corners,align=center,text width=57mm,minimum height=12mm}]
\node[b] (a) at (0,0) {One balanced sextuple\\at its primitive conductor};
\node[b] (l) at (-3.35,-1.8) {Positive predicates\\$\mathrm{Dec}\cup\mathrm{Std}_5$};
\node[b] (r) at (3.35,-1.8) {Direct and lowering moves\\of \Cref{def:move}};
\node[b] (ll) at (-3.35,-3.5) {Li's residual list\\six orbits at $33,39,45,105$};
\node[b] (rr) at (3.35,-3.5) {Our terminal list\\$O_{21},O_{33},O_{39}$};
\draw[->] (a) -- (l);\draw[->] (a) -- (r);
\draw[->] (l) -- node[right,font=\scriptsize]{if neither applies} (ll);
\draw[->] (r) -- node[left,font=\scriptsize]{if no move applies} (rr);
\end{tikzpicture}
\caption{Different tests give different exceptional lists. The left list
is the classification stated in \cite{Li26b}; the right is the
profile-descent assertion of \Cref{thm:moves}. The finite comparisons
below verify their compatibility on the listed examples, rather than
establishing either infinite classification by computation.}
\label{fig:li-comparison}
\end{figure}

\subsection{Explicit comparisons at the seed levels}
The six residual representatives in \cite[(3.1)]{Li26b} and our
degree-$21$ terminal representative are displayed below. The degree-$33$
and degree-$39$ constructions, and the two-pair presentations at $45$
and $105$, are already in Jumagulov \cite[Theorems~1.2--1.4]{Jum26};
they are not new constructions of this appendix.

\begin{center}\small
\begin{tabular}{rll}
\toprule
$m$ & Sextuple & Behaviour under our moves\\
\midrule
21 & $(1,4,9,15,16,18)$ & terminal; lies in $\mathrm{Dec}$\\
33 & $(1,4,16,22,25,31)$ & terminal $O_{33}$\\
39 & $(1,7,16,22,34,37)$ & terminal $O_{39}$\\
39 & $(1,14,16,22,29,35)$ & lowering move $(3,1)$\\
45 & $(1,19,20,28,30,37)$ & direct move $(5,1)$\\
105 & $(1,22,43,64,90,95)$ & direct move $(5,1)$\\
105 & $(3,24,50,66,85,87)$ & direct move $(5,3)$\\
\bottomrule
\end{tabular}
\end{center}

At level $21$, partition $O_{21}$ into $(1,4,18)$ and $(9,15,16)$.
Their sums are $2$ and $-2$ modulo $21$, and the completed quadruples
are $(1,4,18,19)$ and $(2,9,15,16)$. The latter is
$\sigma_{3,2}$, and the former is balanced by the identity in
\Cref{lem:exceptional}. This gives a surface contraction, although
$O_{21}$ has no direct or lowering move. Thus a terminal orbit of our
descent need not be residual for Li's positive predicates.

For the second level-$39$ representative $A$, the move $(3,1)$ gives
\[
 A'=(3,12,16,22,29,35),\qquad
 A'+\sigma_{3,1}=A+\{27,12,3,36\}.
\]
The padding is a union of two inverse pairs. Its order profile drops
from $(39,39,39,39,39,39)$ to $(39,39,39,39,13,13)$.
There is also a direct curve-product description:
$A=(1,16,22)+(14,29,35)$, with both triple sums zero modulo $39$.
As in the curve-product criterion, the simultaneous unit orbit in the
tensor of the two curve $H^1$'s is of type $(1,1)$; Lefschetz $(1,1)$
and the algebraic join give its algebraicity. The first degree-$39$
representative has the analogous split $(1,16,22)+(7,34,37)$.
This is the split construction credited to \cite{Jum26}, also used in
\cite[\S3.2]{Li26b}.

For the remaining three rows, the direct moves give the identities
\begin{align*}
 &(1,19,20,28,30,37)+\{5,40,10,35\}\\
 &\hspace{12mm}=\sigma_{5,1}^{(45)}+(5,20,30,35),\\
 &(1,22,43,64,90,95)+\{85,20,100,5\}\\
 &\hspace{12mm}=\sigma_{5,1}^{(105)}+(5,20,90,95),\\
 &(3,24,50,66,85,87)+\{15,90,45,60\}\\
 &\hspace{12mm}=\sigma_{5,3}^{(105)}+(15,50,60,85).
\end{align*}
Here the superscript records the level. Each right-hand quadruple is
balanced, and each padding consists of inverse pairs. Aoki's standard
cycles, Lefschetz $(1,1)$, the join and pair cancellation therefore make
the three sextuples algebraic, by \Cref{lem:algtools,lem:move}.
These are precisely the already attributed presentations of
\cite{Jum26} and \cite[(3.4)--(3.5)]{Li26b}.

The degree-$33$ row is the same $O_{33}$ treated in
\Cref{lem:exceptional}. Its level-$66$ construction is also
\cite[Theorem~1.4]{Jum26}. Hence four of Li's six residual rows have
moves for our descent and two are its terminal orbits; conversely our
third terminal orbit, at $21$, satisfies a positive surface contraction.
There is no discrepancy between the cardinalities three and six.

\subsection{The scope of the finite verification}
The independent standard-library program
\texttt{anc/check-li-comparison.py} checks these seven tuples against
every unit at their displayed levels. It tests all ten partitions into
two triples, every admissible exact standard tuple, and every move for
each of these particular inputs; it verifies the displayed multiset
identities and profiles and checks inflation at the next two odd
multiples $3m$ and $5m$. Its output is
\texttt{anc/li-comparison-checks.json}.
It is a verification of the named tuples, not an exhaustive census of
all balanced sextuples at these degrees or a proof for unbounded degree.
The infinite classification in \cite{Li26b} rests on its conductor
arguments, not on these checks. The proof of \Cref{thm:oddfour} here
continues to use \Cref{thm:sextuples,thm:moves,lem:exceptional}; it does
not replace any of those arguments by a finite computation or by a
citation to Li.
